% Figure: the Clausius-Clapeyron plot. Natural log of saturation vapour
% pressure against inverse temperature for water, a near-straight line whose
% slope is minus the latent heat of vaporization over the gas constant. The
% straightness of ln p vs 1/T is the working content of the integrated
% Clausius-Clapeyron relation for a constant latent heat.
\begin{figure}[htbp]
\centering
\begin{tikzpicture}
\begin{axis}[
  width=12cm, height=7cm,
  xlabel={inverse temperature $1000/T$ ($\si{\kelvin}^{-1}$)},
  ylabel={$\ln(p_{\mathrm{sat}}/\text{Pa})$},
  xmin=2.6, xmax=3.7, ymin=6, ymax=12,
  axis lines=left,
  legend pos=south west, legend cell align=left,
  legend style={font=\scriptsize, draw=enggray!40},
  tick label style={font=\scriptsize},
  label style={font=\small},
]
% integrated Clausius-Clapeyron: ln p = ln A - L/(R T)
% water: L = 40.7 kJ/mol, R = 8.314; L/R = 4895 K.
% ln p [Pa] = C - 4895/T ; anchor at T=373.15, p=101325 => ln p = 11.526
% => C = 11.526 + 4895/373.15 = 11.526 + 13.119 = 24.645
% plot vs u = 1000/T, so T = 1000/u, ln p = 24.645 - 4.895*u
\addplot[engblue, very thick, domain=2.6:3.7, samples=2]
  {24.645 - 4.895*x};
\addlegendentry{$\ln p = C - L/(RT)$, $L=40.7\,\text{kJ}/\text{mol}$}
% mark the normal boiling point at T=373.15 => 1000/T = 2.680, ln p = 11.526
\addplot[engred, only marks, mark=*, mark size=2pt] coordinates {(2.680,11.526)};
\node[engred, font=\scriptsize, anchor=south east] at (axis cs:2.75,11.4)
  {boiling point, $101\,\text{kPa}$};
% mark 273.15 (freezing) => 1000/T = 3.661, ln p = 24.645 - 17.92 = 6.72
\addplot[engamber, only marks, mark=*, mark size=2pt] coordinates {(3.661,6.727)};
\node[engamber, font=\scriptsize, anchor=north west] at (axis cs:3.4,7.3)
  {$0\,\degC$, $\approx0.6\,\text{kPa}$};
\end{axis}
\end{tikzpicture}
\caption[Clausius-Clapeyron plot of water's vapour pressure]{The logarithm of
water's saturation vapour pressure against inverse temperature. Integrating the
Clausius-Clapeyron relation $d\ln p/dT = L/(RT^2)$ with a latent heat $L$ taken
constant gives $\ln p = C - L/(RT)$, a straight line in these axes whose slope
is $-L/R$. The two marked points are the normal boiling point, where the vapour
pressure reaches one atmosphere, and the triple-point region near $0\,\degC$,
where it falls to under a hundredth of that. The near-straightness of real
vapour-pressure data on this plot is what lets an engineer read a latent heat
off two pressure measurements at two temperatures, and it is the second law at
work: the slope encodes the entropy of vaporization $L/T$, the entropy the
liquid gains in becoming vapour, which the phase equilibrium holds in balance
against the pressure difference across the coexistence line.}
\label{fig:vol04ch04:clausius-clapeyron}
\end{figure}
